La seguridad en entornos habitacionales e industriales representa uno de los pilares fundamentales del bienestar humano y la preservación de activos materiales. En particular, la prevención y detección temprana de incendios y fugas de gases peligrosos constituye un desafío técnico crítico, dado que el tiempo de respuesta es inversamente proporcional a la magnitud de los daños resultantes. Históricamente, la ingeniería ha abordado esta problemática mediante dispositivos aislados; sin embargo, la complejidad de las infraestructuras modernas demanda soluciones integrales que operen de manera coordinada.

En la última década, la arquitectura de la \textit{Internet de las Cosas} (IoT) ha transformado radicalmente este panorama. El despliegue de redes de sensores inalámbricos (\textit{Wireless Sensor Networks}) permite no solo capturar variables físicas como temperatura, humedad y concentración de gases combustibles, sino también procesar dicha información de forma distribuida. Según \citet{valvano2014volume1}, un sistema embebido moderno no debe limitarse a la captura de datos, sino que debe integrar capacidades de comunicación y procesamiento en tiempo real para tomar decisiones autónomas y eficientes.

Bajo esta premisa, el proyecto \textit{NodeAlert IoT } evoluciona de un prototipo de monitoreo único a un ecosistema distribuido basado en la arquitectura \textit{Edge Computing}. Esta transición permite que el procesamiento de señales críticas basado en los principios fundamentales de circuitos eléctricos \citep{alexander2006}  se realice en la periferia de la red (nodos ESP32), reduciendo la latencia y la dependencia de un servidor central. 

La implementación de este sistema se fundamenta en la integración de protocolos de comunicación robustos como MQTT y el uso de sistemas operativos de tiempo real (RTOS) para la gestión de tareas concurrentes . A lo largo de este documento, se detallará cómo esta arquitectura no solo mejora la confiabilidad de la detección, sino que también ofrece una escalabilidad industrial, permitiendo la supervisión remota mediante interfaces web, consolidando una solución accesible y de alto rendimiento para la gestión de emergencias ambientales.